\begin{problem}[Continuous functions really form a sheaf]
Let $X$ be any topological space and $Y$ any topological space. Prove that
\[
\mathcal C_Y(U)=C^0(U,Y)
\]
is a sheaf of sets. Then specialize to $Y=\RR$ and prove that the ring operations make $\mathcal C_\RR$ a sheaf of rings.
\end{problem}

\paragraph{Why this problem matters.}
This is the archetype: a property that is local on the domain gives a sheaf.

\paragraph{Useful ideas before starting.}
Compatible maps glue pointwise. Continuity can be checked on an open cover.

\paragraph{Guided hints.}
For continuity of the glued map $f$, compute $f^{-1}(V)$ for an open $V\subseteq Y$ as a union of its intersections with the cover members.

\ifdefined\FullProblemDossiers
\begin{solution}
Given compatible $f_i:U_i\to Y$, define $f(x)=f_i(x)$ when $x\in U_i$. Compatibility makes this well defined and uniqueness is pointwise. For open $V\subseteq Y$, $f^{-1}(V)\cap U_i=f_i^{-1}(V)$ is open in $U_i$, hence $f^{-1}(V)$ is open because the $U_i$ cover $U$. Thus $f$ is continuous. Locality is immediate. For $Y=\RR$, addition and multiplication are defined pointwise and commute with restriction, so all restriction maps are ring homomorphisms.
\end{solution}

\paragraph{What happened mathematically.}
The local character of continuity is exactly what gluing needs.

\paragraph{Extensions.}
Repeat for smooth, holomorphic, and regular maps in contexts where those properties are local.

\paragraph{Provenance.}
Legacy-derived from the sheaf examples in AG2 and the preferred sheaf notes; expanded as a canonical dossier.
\else
\paragraph{Provenance.}
Legacy-derived from the sheaf examples in AG2 and the preferred sheaf notes; expanded as a canonical dossier.
\fi
